%%%%%%%%%%%%%%%%%%%%%%%%%%%%%%%%%%%%%%%%%%%%%%%%%%%%%%%%%%%%%%%%%%%%%%%%
\section{The continuum model}\label{sec:model}
%%%%%%%%%%%%%%%%%%%%%%%%%%%%%%%%%%%%%%%%%%%%%%%%%%%%%%%%%%%%%%%%%%%%%%%%

This section describes the model of Ang, Park, Pfeffer and Sheffield~\cite[Section~6]{APPS}.
Nothing in it is new.

\subsection{The Gaussian free field}

The \emph{Gaussian free field} is the two-dimensional analog of Brownian motion: a centered
Gaussian random function $h$ of a point $z$ of the plane or the sphere, whose values at two points
have covariance about
\[
\E[h(z)h(w)]\;\approx\;-\log|z-w| .
\]
This is infinite when $z=w$, so $h$ has no values at points; it is a random generalized function.
Its averages over regions are ordinary Gaussian variables. For a square $S$ with side length $|S|$,
the average $h_S$ of $h$ over $S$ is centered normal with
\[
\operatorname{Var}(h_S)=\log\frac1{|S|}+O(1),
\]
and for nested squares $S'\subset S$ the covariance of $h_{S'}$ and $h_S$ is
$\log(1/|S|)+O(1)$. Intuitively, along nested squares each halving of the side adds variance
$\log2$, nearly independently of what came before, so the averages behave like a Brownian motion
in the variable $\log(1/|S|)$. Averages over squares at distance $r$ are correlated by about
$\log(1/r)$.

Formally, $h$ has a density proportional to $\exp\bigl(-\tfrac1{4\pi}\int|\nabla h|^2\bigr)$, where the
integral is the \emph{Dirichlet energy}. Concretely, if $f_1,f_2,\ldots$ are orthonormal for the inner
product $(f,g)_\nabla=\frac1{2\pi}\int\nabla f\cdot\nabla g$ and $a_1,a_2,\ldots$ are independent
standard normal variables, then $h=\sum_ia_if_i$, a sum that converges as a generalized function.
Intuitively, $h$ is a standard Gaussian vector in infinitely many coordinates. This description
makes a determinant computation exact (Section~\ref{sec:tilt}).

APPS work on the sphere with a smooth Riemannian metric $g_0$, a smooth way of measuring lengths
and areas, chosen flat near a closed unit square $S_0$ drawn on the sphere. On the sphere the field
is defined only up to an additive constant; APPS fix the constant by requiring that the integral of
$h$ against the Gaussian curvature of $g_0$ vanish. Near $S_0$ the covariance of $h$ is then $-\log|z-w|$ plus a bounded
smooth function. All estimates below use this exact covariance; none replaces the field by a model
with independent pieces.

\subsection{The dyadic subdivision}

A \emph{dyadic square} is the unit square $S_0$ or one of the squares obtained from it by
repeatedly cutting a square into four equal squares. A dyadic square of \emph{depth} $j$ has side
$2^{-j}$; its \emph{ancestors} are the dyadic squares that contain it.

Fix $Q\ge2$ and a threshold $\varepsilon>0$. Call a dyadic square $S$ \emph{active} if
\begin{equation}\label{eq:rule}
e^{h_S/Q}\,|S|\;>\;\varepsilon .
\end{equation}
Start from $S_0$; cut every active square into four, and repeat on the new squares. The squares
that are never cut form a partition $\Pt^\varepsilon$ of $S_0$, and $N$ is their number. Each cut
adds three squares, so $N\equiv1\pmod 3$. For $Q\ge2$ the procedure stops almost surely, so $N$ is
finite~\cite{APPS}; at $Q=2$ it also follows from $\E[N]<\infty$ (Proposition~\ref{prop:critupper}). Figure~\ref{fig:subdivision} shows a schematic example.

\begin{figure}[htb]
\centering
\begin{tikzpicture}[scale=3.0]
  \draw[line width=0.8pt] (0,0) rectangle (1,1);
  \draw (0.5,0)--(0.5,1); \draw (0,0.5)--(1,0.5);
  % refine the upper-left quarter
  \draw (0.25,0.5)--(0.25,1); \draw (0,0.75)--(0.5,0.75);
  % refine one of its children twice
  \draw (0.375,0.75)--(0.375,1); \draw (0.25,0.875)--(0.5,0.875);
  \draw (0.4375,0.875)--(0.4375,1); \draw (0.375,0.9375)--(0.5,0.9375);
  % refine another child once
  \draw (0.125,0.5)--(0.125,0.75); \draw (0,0.625)--(0.25,0.625);
  % refine the lower-right quarter once
  \draw (0.75,0)--(0.75,0.5); \draw (0.5,0.25)--(1,0.25);
  \node[font=\small] at (0.75,0.75) {low};
  \node[font=\small,anchor=west] at (1.05,0.92) {high field: many small squares};
  \draw[->,>=Stealth] (1.04,0.92)--(0.47,0.93);
\end{tikzpicture}
\caption{A schematic dyadic subdivision. A square is cut while $e^{h_S/Q}|S|>\varepsilon$, so the
subdivision is finest where the field is high. Here $N=19$.}
\label{fig:subdivision}
\end{figure}

Intuitively, $e^{h_S/Q}|S|$ is the size of $S$ as seen by the field: where $h$ is high, squares
count as large and are cut again and again. The final squares play the role of the faces of a
random planar map (Section~\ref{sec:intro}), and $N$ is the size of the surface. The parameter $Q$ enters only through the
rule~\eqref{eq:rule}. Duplantier and Sheffield used the same subdivisions in their proof of a
version of the KPZ relation~\cite{DS}. The \emph{height} of $S$,
\begin{equation}\label{eq:height}
H(S)\;=\;h_S-Q\log\frac{\varepsilon}{|S|},
\end{equation}
measures how far $S$ is above the threshold: $S$ is active exactly when $H(S)>0$, and at depth $j$,
$H(S)=h_S-Qj\log2-Q\log\varepsilon$. From a square to a child the height changes by $-Q\log2$ plus
a random amount of variance about $\log2$. So the squares form a \emph{branching random walk}, a
population in which each individual (a square) has four children whose positions (heights) are
its own plus a random step. Here the steps have negative drift, the population is stopped where
the height becomes negative, and the branches are correlated because nearby squares see nearly the
same field. Section~\ref{sec:brw} studies the version without correlations.

\subsection{The parameters}

APPS tie $Q$ to the central charge $c$ by $Q^2=(25-c)/6$, so $Q\ge2$ corresponds to $c\le1$. The
probability literature writes $Q$ through a second number $\gamma\in(0,2]$, the parameter of
Liouville quantum gravity (Section~\ref{sec:upper}):
\begin{equation}\label{eq:Q}
Q^2=\frac{25-c}{6},\qquad Q=\frac2\gamma+\frac\gamma2,\qquad \gamma=Q-\sqrt{Q^2-4}.
\end{equation}
So $\gamma=2$ at $Q=2$. Squaring the second relation gives
$Q^2-2=4/\gamma^2+\gamma^2/4$. With $x=4/\gamma^2$ this reads
\begin{equation}\label{eq:xQ}
x+\frac1x\;=\;Q^2-2\;=\;\frac{13-c}{6},
\end{equation}
the KPZ relation~\eqref{eq:kpz}, with $x\ge1$ because $\gamma\le2$. So the KPZ formula says that
the tail exponent of the size equals $4/\gamma^2$. The recurring values are: for a bare surface,
$c=0$, $Q^2=25/6$, $\gamma^2=8/3$, $x=3/2$; for one spanning tree, $c=-2$, $Q^2=9/2$, $\gamma^2=2$,
$x=2$; for two curves as in a meander, $c=-4$, $Q^2=29/6$, $\gamma^2=(17-\sqrt{145})/3\approx1.6528$,
$x=\xstar$; and at the critical endpoint, $c=1$, $Q=\gamma=2$, $x=1$.

